%% text-growth.tex -- the growth stage (used by florilegium-grow.tex only)
\subsubsection{The growth stage}

This document is the answer to the exercise, and it goes one step further. Its margin is not
the authored composition of the previous section: nothing in it was placed by hand. A Lua
program decides where every stem, leaf and flower goes and how large it is, and before that,
what kind of plant each stem is.

\begin{definition}
A \emph{genome} is a sequence of traits $\tau_1,\dots,\tau_{62}$, each with a distribution over
its values that may depend on the values already drawn. A \emph{species} is one draw of every
trait, in that order, from a seeded generator. The first trait is the \emph{body plan} (lily,
iris, rose, mustard, pea, mint, daisy, umbellifer, heath, borage, catkin, columbine, a generic
eudicot or monocot, or a fern), which fixes the clade and chooses the distributions of all the
others. Then come the visitor the flowers are for, the flower (its color, form and symmetry,
the number of its parts, the petal's shape, tip, edge and markings, how it is held), the
cluster, the leaf (simple or compound, its shape, base, tip, edge, veins and arrangement), the
stem, the fruit, and how early in the year it flowers.
\end{definition}

The primitives of the previous section become \emph{organs}. A petal is painted as before, but
how many petals a flower has, whether they are free or fused into a bell, a tube or a pea
flower, and how the flowers stand on their stalks are traits of the species. Because every
distribution is conditioned on the body plan, the rules of botany hold by construction. A lily
has its parts in threes, a pea flower is two-sided with five petals and a pod, a catkin has no
petals and no visitor, and only a radially symmetric flower varies its number of petals. Each
species can also describe itself, as ``a blue wheel of five lobes in a coiled cyme that
uncurls, for bees; hairy, smooth-edged, oval leaves; then nutlets''.

\begin{construction}[Community]
The document's seed draws a \emph{pool} of eight species, two heroes, two sprays, a monocot
and three others, with no body plan twice, and with heroes that are \emph{showy} (single
flowers or loose clusters, painted in the ink). Each page draws four of the eight, a hero, a
second showy plant, the monocot and a spray. The same plants therefore recur through the
document, as plants recur through a garden, while each page's selection and layout are its
own.
\end{construction}

\begin{construction}[Stems]
The root is placed near the bottom-left corner of the page. Three stems rise from it, one
plant each; the $x$ coordinate of each is a smooth function of the height,
\[
  x(y)=c+A\,n\bigl(y/\lambda+\varphi\bigr),
  \qquad A=0.2\,\text{(column width)},\ \lambda\in[6,9]\,\text{cm},
\]
with $n$ a two-octave gradient noise and $\varphi$ a phase drawn per stem, blended into the
root over the first four centimeters. Different phases make the stems sway apart and cross
now and then, as real stems do. The hero's stem reaches the top of the page and ends in its
largest flower; the second showy plant stops at $60$ to $72\%$ of the height, the monocot at
$36$ to $48\%$, each ending in its own cluster. A flower seen from the side, such as a bell,
a tube or a pea flower, reaches out from its stalk as a solid shape, so it is drawn at $0.72$
of the size of one seen face on. The spray leaves the root along the bottom of
the page as the footer. Leaves are hung at the nodes as the species sets them, alone, in pairs
or in whorls, $35^\circ$ to $60^\circ$ off the stem, $2.4\,$cm long at the bottom of the page
and smaller toward the top by the species' own taper, and only where there is room; most
monocots raise their strap leaves from the root instead, the longer ones arching over and
some twisting to show their paler underside.
\end{construction}

\begin{definition}
A \emph{proxy} for a motif is a disc of radius $r(s)$, where $s$ is the motif's scale, that
approximates the area it occupies. A point $p$ of the free region $R$ has \emph{free
radius}
\[
  f(p)=\min\Bigl(\,d(p,\partial R),\ \min_i\bigl(\lVert p-c_i\rVert-\kappa\,r_i\bigr)\Bigr),
\]
the radius of the largest disc about $p$ that stays inside $R$ and overlaps no proxy
$(c_i,r_i)$ by more than the fraction $1-\kappa$. With $\kappa=0.85$, a leaf may tuck under
the edge of a flower; stems are proxies too, as chains of small discs.
\end{definition}

\begin{construction}[Growth]
The region $R$ is the margin column, a band along the bottom of the page, and on the last
page also the space between the end of the text and the page number. On a grid of candidates
spaced $2\,$mm apart, each point keeps $f(p)$ and its distance to the nearest stem. Then,
until the largest free radius falls below $3.6\,$mm:
\begin{enumerate}
  \item take the point $p$ that maximizes $f(p)$ among those within $3\,$cm of a stem, with
        a random factor in $[0.85,1.15]$ so that ties break differently on every seed;
  \item choose, by weight, a motif that fits from those of the plant whose stem is nearest,
        which are its cluster, a smaller spray of it, a spray of buds, a leaf, a tendril if it
        climbs, and the visitor its flowers are for. A motif that has already appeared $k$ times has its
        weight divided by $1+k/c$, and the full cluster has a hard cap. Set its scale so that
        $r(s)\approx0.95\,f(p)$;
  \item give it an axis pointing away from the nearest stem point, leaning upward for the
        kinds that reach for the light, and grow a branch to it from that point, as in the
        authored composition. The branch becomes a stem of the same plant for later branches.
\end{enumerate}
Because the largest hole is always taken first, the big clusters are placed early, sprays and
leaves fill the middle, and a last pass puts buds and small leaves on the stems wherever a hole
of more than $1.7\,$mm is left next to one; a hole far from any stem may get one of six copper
accents, and otherwise stays empty. That ordering is the whole reason the result looks
\emph{composed} rather than scattered.
\end{construction}

\begin{remark}[Adapting to the text]
The free region is computed from the page layout, so changing the margins moves the
garden. The end of the text is measured where it happens: \verb|\florilegiumtextend|
records \verb|\pagetotal| and \verb|\pagedepth|, the height of the material already on the
current page, and the background of that page is drawn after it. On the last page, the
garden therefore spills into the space below the words in a single compilation: branches
arch out of the main stems into that space, from each plant in turn, and each ends in a
flower.
\end{remark}

\begin{remark}[Visitors, marks and seasons]
Visitors are sampled from genomes of their own, of fifteen kinds, from butterflies, hawkmoths and
bumblebees to dragonflies, ladybirds, caterpillars, snails and birds, each with its own markings
and pose. A plant brings the visitor its flowers are for, which takes a hole near that plant's
stems like a motif but needs no branch; a snail climbs a main stem, a seed-eater perches on a
twig of its own, a caterpillar feeds on a leaf it has bitten, a dragonfly hunts by the monocot.
Their decisions come from a random stream of their own, so a page without them is the page it
would have been. A mark in the text, left by \verb|\florilegiumbloom| or by a section heading, is
found in the page box at shipout and gets the flowers of the nearest plant, in copper, before any
leaf is hung. An age between $0$ and $1$ sets each flower's stage, from bud to fruit, by the
season, its place in its cluster and how early its species flowers.
\end{remark}

\begin{remark}[The same species in \TeX]
The species are drawn from pgf's own random generator, which the Lua replays exactly, and
every choice is a whole-number draw, so the TikZ library of this package finds the same species
for the same seed with pgfmath alone, under pdf\LaTeX. It grows and paints them with the same
steps and draws, in pgfmath's arithmetic instead of floating point, and the two drawings agree
within a fraction of a point. The layout of a margin around a page's text stays in Lua.
\end{remark}

\begin{remark}[Cost]
The layout is cheap: the grid has a few thousand points, a proxy only updates the points
near it, and the whole loop takes a fraction of a second. Painting is most of the cost. A page
of species is about five megabytes of path operators, one in the PDF, and adds a little over a
second to the compilation.
\end{remark}
